% P16-C4: all ID branches, positive-law invariants, termination and soundness.
\section{The recursive identification algorithm: termination and soundness}
\label{sec:p16-identification-algorithm}

This treatment gives every branch of the algorithm in Shpitser and
Pearl's Figure~3 and proves termination and the correctness of every
successful return \cite[Figure~3, Lemmas~3--8 and Theorem~5,
pp.~1223--1225]{ShpitserPearl2006JointIdentification}.
The execution theorem below has an explicit strictly positive observed-law
domain. The separate do-calculus derivations of the source's utility
lemmas are supplied as well as the preceding latent-sum foundations.
The hedge construction needed for a general completeness conclusion is
a distinct next result.

\begin{hbdefinition}{def:p16-id-call-data}{Global arguments and the invariant of a recursive call}
Fix the observed set \(N\), finite nonempty value sets \(S_i\),
and the independent-latent-input model class of
\cref{def:p16-latent-input-model}. The input graph \(G\) is
supplied as a causal assumption. The given observed law \(p\) is
compatible with this class and satisfies
\[
 p(v)>0\qquad\text{for every }v\in S_N.
\]
The graph is not inferred from positivity or from the observed table.

A call \(\mathsf{ID}(Y,X,q_R,G[R])\) has current random
vertices \(R\subseteq N\), fixed vertices \(F=N\setminus R\),
and disjoint action/response sets \(X,Y\subseteq R\).
Its kernel \(q_R(v_R;v_F)\) is represented as a function of
the full assignment \(v\). The semantic invariant is
\(q_R(v)=Q[R](v)\) for every compatible original model, and
\(q_R(v)>0\) at every full assignment. Each current model is
the proved restriction after fixing \(F\), retaining its actual
latent inputs and their child sets intersected with \(R\).
Marginal or conditional notation for \(q_R\) sums over its other
random coordinates in \(R\), while keeping \(F\) fixed.

For a function \(f:S_N\to\mathbb R\) and \(B\subseteq N\),
write
\[
 (\Sigma_B f)(v)=\sum_{b\in S_B}f(v[B\leftarrow b]),
\]
where \(v[B\leftarrow b]\) replaces exactly the indicated
coordinates. This notation binds summation coordinates explicitly.
The call's target function is
\[
 \mathcal T(R;Y,X)(v)=
      \bigl(\Sigma_{R\setminus(X\cup Y)}Q[R\setminus X]\bigr)(v).
\]
Its free parameters are the response assignment \(v_Y\), action
assignment \(v_X\), and current fixed assignment \(v_F\).
At the initial call \(R=N,q_R=p\), it is the requested effect
\(P_x(y)\). The source passes the assignments \(y,x\) directly;
here they remain explicit arguments of the returned function.

Choose a topological order of the directed part of \(G\) and use
its restriction in each call. For a maximal district \(D\) of
\(G[R]\), let \(\operatorname{Ext}^R_D(q_R)\) be the product
of the current prefix ratios at the vertices of \(D\), as defined
in \cref{thm:p16-latent-district-recovery}. Values outside \(R\)
are held fixed in every prefix sum. Extracting to \(D\) makes
\(N\setminus D\) the new fixed set; those parameter values are
retained, not conditioned away.
\end{hbdefinition}

\begin{hblemma}{thm:p16-id-positive-support}{proved=annex:p16.id-positive-support}{The support required by every recursive operation}
In a current model with positive law \(q_R\), every additional
fixed intervention has positive mass at every assignment of its
remaining random coordinates. Ancestral marginalization preserves
strict positivity. District extraction from a positive kernel is
defined at every assignment, is strictly positive, and equals the
corresponding normalized \(Q\) kernel. Consequently every finite
conditional expression on observed coordinates used below has a positive denominator when
its observed and intervened index sets are disjoint.
\end{hblemma}

\subsection{The source's do-calculus utility derivations}

For this subsection let \(\Gamma=G[R]\), and let
\(\widehat\Gamma\) be its proved expansion with the retained
independent latent roots. Write \(P^R_x(y\mid w)\) for the
conditional intervention law in this current model, with the global
fixed parameters suppressed. Write
\(\mathcal Q_R[D](v)=P^R_{v_{R\setminus D}}(v_D)\). By the restriction
bridge and combination of disjoint interventions, this is \(Q[D](v)\).
All graphical licenses below are statements in the expanded DAG.

\begin{hblemma}{thm:p16-id-do-licenses}{proved=annex:p16.id-do-licenses}{Graphical licenses in the ordered latent-root model}
The following licenses hold in a current call.
\begin{enumerate}
\item For pairwise disjoint observed sets \(X,Y,Z,W\), if no
      member of \(Z\) is an ancestor of \(Y\cup W\) after
      deleting incoming arrows at \(X\), Rule 3 permits insertion
      or deletion of its action given \(W\).
\item Let \(D\) be a maximal district, \(i\in D\), and
      \(L_i\) the observed vertices preceding \(i\). Put
      \(A_D=\operatorname{an}_{\Gamma}(D)\setminus D\),
      \(B_i=A_D\cap L_i\), \(D_i=D\cap L_i\), and
      \(E_i=L_i\setminus(D\cup A_D)\), and let
      \(b_i=v_{B_i}\). In
      \(\widehat\Gamma_{\overline{B_i}}\), the vertices
      \(i,E_i\) are separated given \(B_i\cup D_i\).
      Rule 1 therefore permits insertion of the observations \(E_i\)
      into \(P^R_{b_i}(v_i\mid v_{D_i})\).
\item For any \(B\subseteq L_i\setminus D\), the vertices
      \(i,B\) are separated given \(L_i\setminus B\) in
      \(\widehat\Gamma_{\underline B}\). Rule 2 therefore
      exchanges those actions with observations.
\end{enumerate}
The positivity lemma supplies the required probabilities for all
these applications. No bound of two children per latent input is used.
\end{hblemma}

\begin{hbtheorem}{thm:p16-id-do-utility}{proved=annex:p16.id-do-utility}{Complete do-calculus derivations for source Lemmas 4--8}
The following identities have finite derivations by the three rules
and ordinary probability operations, on the stated positive domain.
\begin{enumerate}
\item For a maximal district \(D\) of \(\Gamma\),
      \(\mathcal Q_R[D]=\operatorname{Ext}^R_D(q_R)\).
      For \(S=R\setminus X\), the general district identity is
      \[
       P^R_x(y)=\Sigma_{S\setminus Y}
               \prod_{D\in\mathcal D_S}\mathcal Q_R[D].
      \]
\item For \(A=\operatorname{an}_{\Gamma}(Y)\) and
      \(q_A=\Sigma_{R\setminus A}q_R\),
      \(P^R_x(y)=P^A_{x_{X\cap A}}(y)\).
\item For \(W=(R\setminus X)\setminus
      \operatorname{an}_{\Gamma[R\setminus X]}(Y)\),
      \(P^R_x(y)=P^R_{x,w}(y)\) for every value assignment \(w\).
\item If \(S=R\setminus X\) is a maximal district of \(\Gamma\),
      \(P^R_x(y)=\Sigma_{S\setminus Y}\operatorname{Ext}^R_S(q_R)\).
\item If \(S\subsetneq D\) for a maximal district \(D\), and
      \(q_D=\operatorname{Ext}^R_D(q_R)\), then
      \(P^R_x(y)=P^D_{x_{X\cap D}}(y)\), with the values on
      \(R\setminus D\subseteq X\) retained as fixed parameters.
\end{enumerate}
These are respectively the district, ancestral, additional-action,
maximal-district base, and larger-district recursion steps in the source.
\end{hbtheorem}

\subsection{All seven recursive branches}

\begin{hbdefinition}{def:p16-id-failure-record}{A retained structural failure record}
A failure contains the raw pair \(G[R],G[S]\) returned at
branch 5, the local sets \((R,X,Y,S)\), the fixed set
\(N\setminus R\), the unchanged observed value sets, and the
sequence of parent call frames and branch numbers. A frame records
its kernel expression and any reference assignment introduced by
branch 3. The raw graph pair is the source's
\(\mathsf{FAIL}(G,S)\) object. It is retained for the subsequent
edge-subgraph hedge construction; it is not silently replaced by
an already proved nonidentification conclusion.
\end{hbdefinition}

\begin{hbmethod}{def:p16-id-algorithm}{The recursive algorithm ID}
Input \((Y,X,q_R,G[R])\) satisfies \cref{def:p16-id-call-data}.
Apply the following branches in order; a return ends the call.
\begin{enumerate}
\item If \(X=\varnothing\), return \(\Sigma_{R\setminus Y}q_R\).
\item Put \(A=\operatorname{an}_{G[R]}(Y)\). If \(A\ne R\),
      return \(\mathsf{ID}(Y,X\cap A,
                  \Sigma_{R\setminus A}q_R,G[A])\).
\item Put \(S=R\setminus X\) and
      \(W=S\setminus\operatorname{an}_{G[S]}(Y)\).
      If \(W\ne\varnothing\), choose any reference tuple
      \(w^\circ\in S_W\), call
      \(\mathsf{ID}(Y,X\cup W,q_R,G[R])\), and return its
      expression with coordinates \(W\) set to \(w^\circ\).
      Retain this choice in the frame. The proved identity makes
      the successful value independent of that choice.
\item Compute \(\mathcal D_S=\{D_1,\ldots,D_k\}\).
      If \(k\geq2\), obtain
      \(e_j=\mathsf{ID}(D_j,R\setminus D_j,q_R,G[R])\)
      for every district, and return
      \(\Sigma_{S\setminus Y}\prod_{j=1}^k e_j\).
\item Otherwise \(\mathcal D_S=\{S\}\). If
      \(\mathcal D_R=\{R\}\), return
      \(\mathsf{FAIL}(G[R],G[S])\) with the specified record.
\item If \(S\in\mathcal D_R\), return
      \(\Sigma_{S\setminus Y}\operatorname{Ext}^R_S(q_R)\).
\item Otherwise take the unique \(D\in\mathcal D_R\) containing
      \(S\), form \(q_D=\operatorname{Ext}^R_D(q_R)\), and
      return \(\mathsf{ID}(Y,X\cap D,q_D,G[D])\).
\end{enumerate}
Any failed child propagates its failure record immediately, with its
parent frame retained; no product, sum or substitution is applied to
a failure. A fixed ordering of districts can be used to determine
which failed child is encountered first. Summation coordinates in
different kernel expressions are renamed or interpreted through
\(\Sigma_B\), so an action argument cannot be accidentally captured.
\end{hbmethod}

\begin{hbtheorem}{thm:p16-id-invariants}{proved=annex:p16.id-invariants}{Every branch is defined and preserves the call invariant}
Starting from \(\mathsf{ID}(Y,X,p,G)\), every reached call
has the stated random/fixed sets, actual kernel \(Q[R]\),
strict positivity, and disjoint action/response sets. The branch-7
district always exists and satisfies \(S\subsetneq D\subsetneq R\).
At branch 5,
\[
 \varnothing\ne Y\subseteq S\subsetneq R,\quad X=R\setminus S,
 \quad R=\operatorname{an}_{G[R]}(Y),\quad
 S=\operatorname{an}_{G[S]}(Y),
\]
and both \(R\) and \(S\) are single maximal districts in
their respective recorded graphs. These are the structural facts
certified by the failure record.
\end{hbtheorem}

\begin{hbtheorem}{thm:p16-id-termination}{proved=annex:p16.id-termination}{Termination of the complete recursion}
Each recursive child strictly decreases
\[
 \nu(R,X)=|R|+|R\setminus X|.
\]
The branching is finite. Every run therefore terminates with an
expression or a failure record, and its call depth is at most
\(2|N|+1\). This is a termination bound, not a claim of polynomial
total running time or numerical bit complexity.
\end{hbtheorem}

\begin{hbtheorem}{thm:p16-id-success-soundness}{proved=annex:p16.id-success-soundness}{Soundness of every successful return}
Every successful call returns exactly \(\mathcal T(R;Y,X)\).
In particular the initial successful expression identifies \(P_x(y)\)
from \(p\) in the stated positive-law model class. It uses finite
sums, products and positive-denominator ratios of the observed law.
Its successful trace also supplies a finite derivation using the
do-calculus utility identities and ordinary probability operations.

For a nonpositive observed law, the preceding foundation's local
support identities remain available, but the execution theorem here
does not certify an unrestricted run. Missing ratio support is a
different issue from the graph failure in branch 5. Neither changes
the existence of the underlying equation-replacement intervention.
\end{hbtheorem}

\subsection{Complete success and failure runs}

\begin{hbdefinition}{def:p16-id-worked-models}{Two strictly positive examples}
For the first example use observed vertices \(A,B,C\) and independent
latent bits \(U,E_A,E_B,E_C\), with \(U\) fair and respective
one-probabilities \(1/5,1/4,1/8\) for the other three. Set
\[
 A=U\oplus E_A,\quad B=A\oplus E_B,\quad
 C=\mathbf1_{\{B=1\text{ or }U=1\}}\oplus E_C.
\]
The observed graph is \(A\to B\to C\), \(A\leftrightarrow C\).
Write \(\alpha(0)=4/5,\alpha(1)=1/5\),
\(\beta(0)=3/4,\beta(1)=1/4\), and
\(\gamma(0)=7/8,\gamma(1)=1/8\).

For the second example use graph \(A\to B\to C\),
\(A\leftrightarrow B\), with independent fair roots \(U,E_B\)
and \(\mathbb P(E_C=1)=1/5\). Set
\[
 A=U,\qquad B=A\oplus U\oplus E_B,\qquad C=B\oplus E_C.
\]
Write \(\eta(0)=4/5,\eta(1)=1/5\).
All displayed equations are total on their binary input sets.
\end{hbdefinition}

\begin{hbtheorem}{thm:p16-id-worked-success}{proved=annex:p16.id-worked-success}{Successful runs through district extraction and additional actions}
The first model has the strictly positive table
\[
\begin{array}{cc|cc}
 a&b&p(a,b,0)&p(a,b,1)\\\hline
 0&0&87/320&33/320\\
 0&1&5/320&35/320\\
 1&0&11/320&29/320\\
 1&1&15/320&105/320
\end{array}
\]
The complete call for \(P_a(c)\) returns
\[
 \sum_b p(b\mid a)\sum_{a'}p(a')p(c\mid a',b).
\]
For \(c=1\), its values are \(19/32\) at \(a=0\) and
\(25/32\) at \(a=1\).

The second model has \(p(a,b,c)=\frac14\eta(c\oplus b)>0\).
For the target \(P_b(c)\), branch 3 adds the action on \(A\),
and the eventual return is \(p(c\mid a^\circ,b)=\eta(c\oplus b)\),
independent of the chosen reference value \(a^\circ\).
The annex gives every call, branch and intermediate kernel in both runs.
\end{hbtheorem}

\begin{hbtheorem}{thm:p16-id-worked-failure}{proved=annex:p16.id-worked-failure}{A complete positive-law failure run}
For the graph \(A\to Y\), \(A\leftrightarrow Y\), and
the uniform observed law \(p(a,y)=1/4\), the call for \(P_a(y)\)
returns the raw pair consisting of this two-vertex graph and the
singleton \(Y\) graph at branch 5. In this particular binary case,
the two positive-law models in
\cref{thm:p16-latent-foundation-limits} independently show
nonidentification, with response-one intervention probabilities
\(3/4\) and \(1/2\).
\end{hbtheorem}

Raw failure graphs need not already be the pruned forests.
\Cref{thm:p16-hedge-failure-extraction} constructs an original-query
hedge from the complete failure trace. The subsequent witness and
lifting proofs give the fixed-domain and support conditions of
\cref{thm:p16-hedge-completeness}, including the required cardinalities
of any vertices carrying both a forest coordinate and a response relay.
